

































\documentclass[11pt,a4paper]{article}

\usepackage[utf8]{inputenc}
\usepackage[T1]{fontenc}
\usepackage[french]{babel}
\usepackage{amsmath,amssymb,amsthm,mathtools}
\usepackage{geometry}
\usepackage{enumitem}
\usepackage{booktabs}
\usepackage{array}
\usepackage{hyperref}
\usepackage{microtype}
\geometry{margin=2.5cm}
\setlist[itemize]{leftmargin=1.7em}
\setlist[enumerate]{leftmargin=1.9em}

\title{\textbf{Infini vs fini -- Formalisation rigoureuse d'une hypothèse}\\
\large Version 1 : contraintes, déterminations et solutions}
\author{}
\date{}

\newtheorem{definition}{Définition}
\newtheorem{proposition}{Proposition}
\newtheorem{remark}{Remarque}
\newtheorem{question}{Question ouverte}

\begin{document}
\maketitle

\begin{abstract}
Ce document formalise une hypothèse métaphysique selon laquelle l'infini est
ontologiquement fondamental, tandis que ce que nous appelons habituellement
« fini » n'est pas une catégorie ontologique primitive mais une détermination
obtenue relativement à des contraintes. La reformulation centrale remplace
la chaîne
\[
\text{réalité non limitée}\to\text{contraintes}\to\text{distinctions}
\to\text{déterminations}\to\text{fini}
\]
par
\[
\boxed{
\text{réalité non limitée}\to\text{contraintes}\to\text{problème}
\to\text{solution}.
}
\]
Le but n'est pas de remplacer cette hypothèse par les définitions standards
des mathématiques, mais de la rendre explicite, cohérente, critiquable et
productrice de conséquences. Les mathématiques, l'algorithmique et les
sciences peuvent ensuite être examinées comme des niveaux particuliers de
cette architecture, sans être utilisées comme tribunal ontologique préalable.
\end{abstract}

\section{Principe de fidélité}

La théorie étudiée ici doit être comprise selon ses propres primitives.
Les termes « infini », « fini », « contrainte », « détermination » et
« solution » ne sont donc pas supposés avoir, dès le départ, exactement le
sens que leur donnent les théories mathématiques ou métaphysiques usuelles.

Cela n'interdit pas la comparaison avec les conceptions standards. Il faut
cependant distinguer trois opérations :
\begin{enumerate}
\item formaliser l'hypothèse telle qu'elle est proposée ;
\item en déduire ses conséquences ;
\item comparer ensuite ces conséquences à d'autres cadres.
\end{enumerate}
Une définition mathématique interne à un système formel ne constitue pas, à
elle seule, une preuve que son objet est une catégorie ontologique absolue.

\section{Thèse fondamentale}

\begin{definition}[Non-limitable]
On appelle ici \emph{infini} \(I\) ce qui, par principe, n'est pas susceptible
d'une limitation, d'une totalisation ou d'une clôture ontologique ultime.
\end{definition}

Cette définition est négative : elle caractérise l'infini par l'absence de
borne ontologique ultime. Elle n'affirme pas que l'infini est simplement
« très grand », ni qu'il est seulement ce qu'un être humain ne peut mesurer.

La thèse métaphysique centrale est :
\[
\boxed{\text{\(I\) est ontologiquement fondamental.}}
\]

Le terme « fondamental » signifie ici : si l'on remonte vers les conditions
ontologiques premières, l'infini ne doit pas être expliqué comme le produit
d'une catégorie plus fondamentale appelée « fini ».

\subsection{Non-totalisabilité}

Dans ce cadre, « totalité de l'infini » est problématique. Si une totalité
est comprise comme une clôture déterminée contenant exhaustivement son objet,
alors appliquer cette notion à \(I\) reviendrait à lui attribuer précisément
une clôture que la définition lui refuse.

Ainsi :
\[
\boxed{\text{totaliser \(I\)}\;\Longrightarrow\;\text{le déterminer comme clos}}
\]
et, dans le cadre présent,
\[
\boxed{\text{\(I\) ne peut pas être ontologiquement clos.}}
\]

Il ne faut cependant pas conclure qu'aucune phrase mathématique contenant
le mot « infini » n'a de sens. La difficulté porte sur l'interprétation
ontologique de la totalité, pas sur la légitimité des formalismes.

\section{Le fini comme détermination}

La conception habituelle peut donner l'impression que « fini » désigne une
propriété première d'une chose. L'hypothèse présente inverse la priorité.

Pour une détermination \(D\), on ne pose pas primitivement :
\[
D\in\text{Fini}.
\]
On demande plutôt :
\[
\boxed{D\text{ est une détermination produite ou sélectionnée sous contraintes}.}
\]

Dans cette perspective, « fini » est une qualification secondaire, dépendant
du cadre de détermination.

Par exemple, le nombre \(123\) peut être traité comme une détermination
numérique sans faire de l'énoncé « \(123\) est fini » une conclusion
ontologique sur ce qu'est fondamentalement \(123\). Ce qui importe est que
\(123\) constitue une valeur déterminée dans un cadre donné.

\section{La structure centrale : contraintes et solutions}

La chaîne fondamentale devient :
\[
\boxed{
I
\overset{C}{\longrightarrow}
P
\overset{S}{\longrightarrow}
D
}
\]
où :
\begin{itemize}
\item \(I\) : réalité supposée non limitée ;
\item \(C\) : ensemble de contraintes ;
\item \(P\) : problème ou espace de détermination imposé par ces contraintes ;
\item \(S\) : processus de résolution ;
\item \(D\) : détermination obtenue.
\end{itemize}

La notion de solution doit être entendue au sens général :
\[
\boxed{
\text{solution}=
\text{détermination compatible avec un ensemble de contraintes}.
}
\]

Elle n'est donc pas réduite à la solution d'une équation.

Une contrainte peut être physique, causale, logique, structurelle,
matérielle, cognitive, computationnelle, conceptuelle ou mathématique.

\subsection{Espace des solutions}

Soit \(C\) un ensemble de contraintes et \(\mathcal D\) un domaine de
déterminations admissibles. On définit :
\[
\mathcal S(C)=
\{D\in\mathcal D\mid D\text{ satisfait }C\}.
\]

Une détermination observée ou construite \(D\) satisfait alors :
\[
D\in\mathcal S(C).
\]

Si l'on ajoute des contraintes, \(C'\supseteq C\), alors, sous l'hypothèse
que les contraintes supplémentaires sont compatibles avec le cadre :
\[
\boxed{\mathcal S(C')\subseteq\mathcal S(C).}
\]

Une contrainte supplémentaire ne crée donc pas nécessairement une nouvelle
réalité fondamentale ; elle peut réduire l'espace des déterminations
admissibles.

\section{Exemple élémentaire}

Pour
\[
x+2=5,
\]
on obtient :
\[
x=3.
\]

Dans le cadre présent, le point conceptuel n'est pas de nier la validité
mathématique de cette solution. Il est de changer l'interprétation
ontologique :
\[
\boxed{3=\text{détermination satisfaisant les contraintes du problème}.}
\]

Le nombre \(3\) est ainsi étudié comme résultat déterminé d'un cadre, et non
comme preuve qu'il existe une catégorie ontologique primitive appelée
« fini ».

\section{Le cas de \(123\)}

La chaîne peut être explicitée ainsi :
\[
C_1\longrightarrow D_1=123.
\]

Un autre système de contraintes peut produire :
\[
C_2\longrightarrow D_2=124.
\]

Le changement de résultat ne signifie pas nécessairement un changement du
fond ontologique \(I\). Il signifie, dans la théorie, que les conditions de
détermination ont changé :
\[
\boxed{
I+C_1\to123,\qquad I+C_2\to124.
}
\]

Le nombre est donc une détermination contextuelle au sens relationnel :
\[
D=D(C).
\]

\section{L'exemple d'un univers contenant 17 particules}

Supposons une description selon laquelle un univers contient exactement
17 particules.

La théorie présente demande de distinguer :
\[
\text{« 17 est la valeur obtenue dans ce cadre »}
\]
de
\[
\text{« 17 constitue une limite ontologique absolue de tout réel possible ».}
\]

La première proposition peut être parfaitement cohérente dans un modèle.
La seconde exige une justification supplémentaire.

On peut représenter :
\[
C_{17}\longrightarrow D_{17}=17.
\]

Dire que les contraintes du modèle fixent 17 ne prouve donc pas, dans cette
métaphysique, que l'infini lui-même est devenu 17 ou qu'une frontière absolue
a été découverte.

La possibilité conceptuelle d'une 18e particule peut être représentée par
une autre détermination compatible avec un cadre modifié :
\[
C_{18}\longrightarrow D_{18}=18.
\]

Il faut toutefois distinguer cette possibilité conceptuelle d'une affirmation
empirique : l'hypothèse métaphysique ne suffit pas, à elle seule, à établir
qu'un univers physique réel possède ou pourrait physiquement posséder cette
particule.

\section{Distinction cruciale : impossibilité de connaître et impossibilité ontologique}

Deux propositions sont logiquement différentes :
\[
\text{« nous ne pouvons pas connaître une limite »}
\]
et
\[
\text{« il n'existe pas de limite ontologique ultime ».}
\]

La première est épistémologique. La seconde est ontologique.

La théorie visée est une thèse du second type. Il serait donc insuffisant de
la défendre seulement par les limites de l'observation humaine.

Réciproquement, l'incapacité humaine à concevoir exhaustivement \(I\) ne
constitue pas à elle seule une démonstration que \(I\) est ontologiquement
non-limitable. Il faut que la non-limitabilité fasse partie du principe
métaphysique de départ.

\section{Statut de la définition humaine de l'infini}

Une définition humaine de \(I\) est une détermination conceptuelle. Elle ne
doit donc pas être identifiée à l'objet métaphysique qu'elle cherche à
désigner.

On peut écrire :
\[
I_{\mathrm{réel}}\neq D_{\mathrm{conceptuelle}}(I)
\]
si le signe \(\neq\) exprime ici la non-identité entre l'objet supposé et sa
représentation conceptuelle, et non une négation de toute connaissance.

La formulation prudente est :
\[
\boxed{
\text{toute représentation humaine de \(I\) est déterminée ;
\(I\) n'est pas épuisé par cette représentation.}
}
\]

\section{Conséquences ontologiques}

\subsection{Renversement de priorité}

La priorité ontologique passe de :
\[
\text{objets finis}\to\text{extension éventuelle vers l'infini}
\]
à :
\[
\boxed{
\text{non-limitable}\to\text{contraintes}\to\text{déterminations}.
}
\]

Le problème philosophique principal n'est alors plus :
\[
\text{« pourquoi existe-t-il du fini ? »}
\]
mais :
\[
\boxed{
\text{« pourquoi et comment existe-t-il des contraintes qui produisent des déterminations ? »}
}
\]

\subsection{Le fini devient relationnel}

Si « fini » est une qualification d'une détermination sous contraintes, alors
sa signification dépend du cadre :
\[
F_C(D).
\]

Une même détermination peut être suffisamment déterminée dans un cadre et
être incluse dans un cadre plus large dans lequel d'autres distinctions
deviennent pertinentes.

Cela ne signifie pas que toute frontière physique est imaginaire. Cela
signifie que la propriété « fini » ne reçoit pas automatiquement un statut
ontologique absolu du seul fait qu'une description possède une borne.

\subsection{Pas d'extérieur ontologique}

Si \(I\) est par définition non-totalisable et non-limitable, parler d'un
« extérieur de \(I\) » comme d'un domaine ontologique séparé soulève une
difficulté : un extérieur qui serait absolument séparé introduirait
précisément une détermination globale ou une frontière.

La conséquence doit cependant être formulée prudemment :
\[
\boxed{
\text{dans ce cadre, aucun « extérieur » ne peut être posé comme limite
ontologique ultime de \(I\) sans modifier la définition de \(I\).}
}
\]

\section{Conséquences épistémologiques et perspectives}

Toute connaissance concrète procède par détermination :
\[
\text{sélection}\to\text{distinction}\to\text{description}.
\]

Le sujet connaissant ne saisit donc jamais l'infini comme un objet fini
épuisable. Il construit des modèles déterminés.

Cela conduit à une conception stratifiée :
\[
I
\to
C_{\text{physiques}}
\to
D_{\text{physiques}}
\to
C_{\text{cognitives}}
\to
D_{\text{conceptuelles}}.
\]

Les modèles ne sont pas nécessairement faux parce qu'ils sont déterminés.
Ils deviennent problématiques seulement lorsqu'une détermination locale est
prise pour une limite ontologique absolue sans argument supplémentaire.

\section{Conséquences mathématiques}

\subsection{Ne pas confondre deux niveaux}

Les mathématiques standard disposent de définitions précises de « fini »,
« infini », ensemble, cardinalité, limite, etc. Ces définitions sont
indispensables à la cohérence des théories concernées.

La présente hypothèse ne les réfute pas. Elle affirme une question différente :
\[
\boxed{
\text{une définition mathématique interne n'est pas automatiquement une
ontologie du réel.}
}
\]

\subsection{Interprétation par contraintes}

Une structure mathématique peut être comprise comme un système de
contraintes. Soit \(C\) un ensemble de conditions et \(X\) un espace de
candidats :
\[
\mathcal S(C)=\{x\in X:C(x)\}.
\]

La notion de solution apparaît alors naturellement :
\[
x\in\mathcal S(C).
\]

Cela permet d'intégrer la perspective proposée sans abolir la distinction
entre syntaxe, sémantique et ontologie.

\subsection{Sous-détermination}

Si :
\[
|\mathcal S(C)|>1,
\]
les contraintes ne déterminent pas une solution unique.

Une détermination supplémentaire \(C'\) peut alors satisfaire :
\[
\mathcal S(C')\subsetneq\mathcal S(C).
\]

L'unicité peut donc être comprise comme un effet de contraintes suffisamment
fortes, plutôt que comme une propriété métaphysique primitive.

\subsection{Attention à une limite logique}

Il ne s'ensuit pas que tout objet mathématique est littéralement une
« solution » d'un problème extérieur à lui. Certaines structures sont
introduites axiomatiquement, et certaines notions ne sont pas des solutions
au sens ordinaire.

La thèse peut être maintenue en parlant plus généralement de
\emph{déterminations sous conditions} ; « solution » est alors le concept
unificateur, pas une réduction de toute mathématique à la résolution
d'équations.

\section{Conséquences algorithmiques et computationnelles}

\subsection{Algorithme comme processus de résolution}

Un algorithme peut être vu comme une procédure :
\[
(C,E)\longrightarrow D,
\]
où \(C\) représente les contraintes et \(E\) les données d'entrée.

Le résultat \(D\) est une détermination produite par la procédure.

Ainsi :
\[
\boxed{
\text{algorithme}=
\text{procédure de recherche ou de construction d'une détermination
satisfaisant des contraintes}.
}
\]

\subsection{Espace de recherche}

Un problème algorithmique peut être modélisé par :
\[
\mathcal S(C,E).
\]

Un algorithme efficace peut réduire progressivement le domaine de recherche :
\[
\mathcal S_0\supseteq\mathcal S_1\supseteq\cdots\supseteq\mathcal S_k.
\]

La résolution correspond alors à l'obtention d'un élément admissible, ou à
la preuve qu'aucun élément admissible n'existe dans le modèle :
\[
D\in\mathcal S_k
\quad\text{ou}\quad
\mathcal S_k=\varnothing.
\]

\subsection{Complexité}

Cette perspective rend visible une distinction importante :
\[
\text{existence d'une solution}\neq\text{accessibilité algorithmique de la solution}.
\]

Une solution peut exister mathématiquement tout en étant coûteuse,
inaccessible dans les ressources disponibles, ou impossible à calculer par
une classe d'algorithmes donnée.

Ainsi, les contraintes comprennent aussi les ressources :
\[
C_{\mathrm{physique}},\quad
C_{\mathrm{mémoire}},\quad
C_{\mathrm{temps}},\quad
C_{\mathrm{calcul}}.
\]

Une solution devient alors une détermination relative non seulement au
problème, mais aussi aux ressources autorisées.

\subsection{Conséquence philosophique}

Cela donne un parallèle structurel :
\[
\boxed{
\text{contraindre l'espace des possibles}
\longleftrightarrow
\text{réduire l'espace de recherche}.
}
\]

Mais l'analogie ne doit pas être transformée en identité : l'univers n'est
pas démontré être un ordinateur, et une ontologie des contraintes n'implique
pas automatiquement une ontologie computationnelle.

\section{Conséquences physiques : portée et prudence}

La théorie suggère une lecture selon laquelle les états physiques déterminés
sont des solutions de contraintes physiques :
\[
C_{\mathrm{phys}}\longrightarrow D_{\mathrm{phys}}.
\]

Cette formulation est compatible avec l'idée générale de lois comme
contraintes, mais elle ne constitue pas encore une théorie physique.

Pour devenir physique, il faudrait spécifier :
\begin{enumerate}
\item quelles sont les contraintes fondamentales ;
\item comment elles sont représentées ;
\item comment elles produisent les états ;
\item quelles observations distinguent cette théorie d'autres théories ;
\item quelles prédictions nouvelles elle produit.
\end{enumerate}

Sans cela, l'hypothèse reste métaphysique.

\section{Conséquences causales et structurelles}

La notion de contrainte permet de distinguer :
\[
\text{cause}\neq\text{contrainte}\neq\text{solution}.
\]

Une cause peut contribuer à la production d'une détermination, tandis qu'une
contrainte définit ce qui est admissible. Plusieurs causes peuvent conduire
à une même solution, et une même contrainte peut admettre plusieurs solutions.

Il faut donc éviter de réduire toute contrainte à une cause.

On peut envisager :
\[
C=\{C_1,C_2,\ldots,C_n\}
\]
et :
\[
D\in\mathcal S(C).
\]

La détermination est alors le résultat d'une compatibilité structurelle,
sans que l'ontologie de la théorie soit obligée de postuler une causalité
simple et linéaire.

\section{Conséquences sur la notion de possibilité}

La théorie distingue naturellement :
\[
\text{possible sous \(C\)}
\quad\Longleftrightarrow\quad
D\in\mathcal S(C).
\]

Une possibilité n'est donc pas nécessairement une réalité actualisée.

On peut avoir :
\[
D\in\mathcal S(C)
\]
sans que \(D\) soit réalisé dans le cadre empirique considéré.

Cette distinction protège la théorie contre une confusion entre :
\[
\text{« déterminable »},\quad
\text{« possible »},\quad
\text{« réalisé »}.
\]

Toutes trois sont des notions différentes, même si elles sont reliées par
la structure des contraintes.

\section{Conséquence sur le langage}

Les mots « fini », « objet », « totalité », « extérieur » et même
« réalité » sont eux-mêmes des déterminations conceptuelles.

Le langage procède donc par découpage :
\[
I\to\text{distinctions linguistiques}\to\text{concepts}.
\]

Le danger est de réifier la découpe :
\[
\text{concept déterminé}
\not\Rightarrow
\text{limite ontologique absolue}.
\]

Cela fournit une règle méthodologique centrale :
\[
\boxed{
\text{ne pas confondre la frontière d'une description avec une frontière de l'être.}
}
\]

\section{Conséquences métaphysiques plus générales}

\subsection{Ontologie des relations plutôt que des bornes}

Si les déterminations sont des solutions sous contraintes, leur identité
ontologique est en partie relationnelle :
\[
D=D(C).
\]

L'objet n'est donc pas supprimé ; son statut est réinterprété. On passe d'une
ontologie où la borne est primitive à une ontologie où la détermination et sa
relation aux contraintes sont primitives.

\subsection{Dynamique des contraintes}

Une théorie plus développée pourrait autoriser :
\[
C_0\to D_0\to C_1\to D_1\to\cdots
\]

Ainsi une solution à un niveau peut devenir une contrainte au niveau suivant.

C'est une conséquence particulièrement féconde :
\[
\boxed{
\text{solution à un niveau}
\;\longrightarrow\;
\text{contrainte au niveau suivant}.
}
\]

Elle permet une architecture récursive sans supposer que l'infini est lui-même
une solution finale.

\subsection{Absence de « solution finale »}

Si \(I\) est non-totalisable, il ne faut pas rechercher une solution \(D^\ast\)
qui clôturerait ontologiquement tout le réel :
\[
D^\ast=\text{solution ultime de tout}.
\]

Une telle solution fonctionnerait comme une totalisation, et entrerait donc en
tension avec la définition initiale de \(I\).

La théorie est ainsi structurellement ouverte :
\[
\boxed{
\text{détermination}\neq\text{clôture ontologique finale}.
}
\]

\section{Ce que la théorie n'établit pas}

Pour rester rigoureux, plusieurs conclusions ne peuvent pas être obtenues
sans prémisses supplémentaires.

Elle n'établit pas, à elle seule :
\begin{itemize}
\item que l'univers physique est réellement infini ;
\item qu'il existe une infinité actuelle d'objets physiques ;
\item que toute réalité est computationnelle ;
\item que toute solution est algorithmique ;
\item que toute frontière physique est fictive ;
\item que toute possibilité conceptuelle est physiquement réalisable ;
\item qu'il existe une loi physique fondamentale particulière ;
\item que les mathématiques standard sont incohérentes.
\end{itemize}

Ces propositions pourraient devenir des hypothèses supplémentaires, mais elles
ne sont pas des conséquences logiques immédiates de la thèse.

\section{Tensions internes à résoudre}

\subsection{D'où viennent les contraintes ?}

La question la plus profonde devient :
\[
\boxed{\text{Pourquoi existe-t-il des contraintes ?}}
\]

Si l'on répond que les contraintes sont elles-mêmes des déterminations,
alors on obtient une régression potentiellement récursive :
\[
I\to C_1\to D_1\to C_2\to D_2\to\cdots.
\]

Cette régression n'est pas nécessairement un défaut, mais il faut décider si
elle est :
\begin{itemize}
\item infinie ;
\item hiérarchique ;
\item circulaire ;
\item ou fondée sur une classe de contraintes primitives.
\end{itemize}

\subsection{Les contraintes sont-elles réelles ou descriptives ?}

Il faut distinguer :
\[
C_{\mathrm{ont}}\quad\text{(contrainte ontologique)}
\]
de
\[
C_{\mathrm{desc}}\quad\text{(contrainte imposée par notre description)}.
\]

Une théorie physique ou métaphysique doit préciser si ses contraintes sont
dans le réel, dans le modèle, ou aux deux niveaux avec des correspondances
explicites.

\subsection{Pourquoi une solution plutôt qu'une autre ?}

Si :
\[
|\mathcal S(C)|>1,
\]
la théorie doit expliquer la sélection d'une détermination particulière :
\[
C\to D_i.
\]

Cette sélection peut demander une contrainte supplémentaire, une dynamique,
une condition initiale ou une notion d'actualisation. Ce point est central
pour passer d'une ontologie des possibilités à une ontologie des faits.

\section{Synthèse formelle}

L'hypothèse peut être condensée en six principes :

\begin{enumerate}
\item \textbf{Fondamentalité :}
\[
I\text{ est ontologiquement premier.}
\]

\item \textbf{Non-limitabilité :}
\[
I\text{ n'admet pas de limite ontologique ultime.}
\]

\item \textbf{Détermination :}
\[
C\to\mathcal S(C).
\]

\item \textbf{Solution :}
\[
D\in\mathcal S(C).
\]

\item \textbf{Relativité du « fini » :}
\[
\text{« fini » décrit une détermination sous un cadre de contraintes,
et n'est pas posé comme primitive ontologique.}
\]

\item \textbf{Ouverture :}
\[
D\neq\text{totalisation finale de }I.
\]
\end{enumerate}

La structure générale devient donc :
\[
\boxed{
I
\longrightarrow
C
\longrightarrow
\mathcal S(C)
\longrightarrow
D
}
\]
avec la possibilité de récursion :
\[
\boxed{
D_n\longrightarrow C_{n+1}\longrightarrow
\mathcal S(C_{n+1})\longrightarrow D_{n+1}.
}
\]

\section{Tableau des conséquences}

\begin{center}
\begin{tabular}{>{\raggedright\arraybackslash}p{0.19\textwidth}
                >{\raggedright\arraybackslash}p{0.34\textwidth}
                >{\raggedright\arraybackslash}p{0.34\textwidth}}
\toprule
Niveau & Réinterprétation & Conséquence principale\\
\midrule
Ontologique & L'infini est fondamental ; le fini n'est pas primitif &
Les bornes déterminées ne sont pas automatiquement absolues.\\
Épistémologique & Toute connaissance est une détermination &
Un modèle limité ne constitue pas nécessairement une limite du réel.\\
Mathématique & Une valeur peut être l'élément d'un espace de solutions &
Les contraintes déterminent des sous-espaces admissibles.\\
Algorithmique & Résoudre = chercher/construire une détermination admissible &
L'espace de recherche est réduit par les contraintes.\\
Physique & Un état peut être modélisé comme solution de contraintes physiques &
Il faut encore spécifier des lois et produire des prédictions.\\
Causal & Contraindre n'est pas identique à causer &
Les contraintes peuvent définir l'admissibilité sans être des causes simples.\\
Modal & Possible sous \(C\) signifie admissible sous \(C\) &
Possible, déterminable et réalisé doivent rester distincts.\\
Linguistique & Les concepts sont des découpages déterminés &
Une frontière conceptuelle ne prouve pas une frontière ontologique.\\
Métaphysique & Les solutions peuvent devenir des contraintes &
La théorie peut former une hiérarchie récursive de déterminations.\\
\bottomrule
\end{tabular}
\end{center}

\section{Conclusion}

La reformulation « le fini est une solution » modifie profondément
l'architecture de l'hypothèse.

La formulation la plus fidèle n'est pas :
\[
\text{« rien n'est fini »},
\]
car cette phrase réintroduirait le fini comme catégorie centrale.

Elle est plutôt :
\[
\boxed{
\text{ce que nous appelons « fini » est une détermination obtenue sous
certaines contraintes ; sa détermination n'implique pas, à elle seule,
une limitation ontologique ultime du fond du réel.}
}
\]

La thèse fondamentale devient alors :
\[
\boxed{
\text{non-limitable}
\rightarrow
\text{contraintes}
\rightarrow
\text{espaces de déterminations}
\rightarrow
\text{solutions}
}
\]

Cette structure déplace la question philosophique. Le problème fondamental
n'est plus d'expliquer comment un « infini » serait construit à partir de
petits objets finis, mais d'expliquer comment des contraintes, des
déterminations et des solutions peuvent émerger, s'organiser et se stabiliser
sans transformer une détermination locale en clôture ontologique de
l'ensemble.

La théorie est donc cohérente comme programme métaphysique à condition de
maintenir ses distinctions : infini versus représentation de l'infini,
contrainte versus cause, solution versus réalisation, possibilité versus
actualité, et formalisation mathématique versus interprétation ontologique.
Ses principales questions ouvertes portent sur l'origine des contraintes,
leur statut ontologique, le mécanisme de sélection des solutions et la
possibilité de dériver de cette ontologie des prédictions qui la distingueraient
empiriquement d'autres cadres.

\end{document}

